\section{Proofs for Section~\ref{sec:parameter-free}}
\label{app:parameter-free-proofs}


\subsection{Uniform-Accuracy Parameter-Free Method}
\label{app:uniform-parameter-free-proofs}


\begin{proof}[Proof of Corollary~\ref{cor:scar-inner}]
For $q_x=-f(x,\cdot)$, strong monotonicity and Lipschitz continuity
of $\nabla q_x$ give
\[
    \mu\|z-y^{\rm in}\|
    \le
    \|\nabla q_x(z)-\nabla q_x(y^{\rm in})\|
    \le
    \ell\|z-y^{\rm in}\|.
\]
Hence
\[
    \mu\le a_x\le\ell.
\]
For $\kappa>1$, the initialization
$\mu_0=M_0=a_x$ therefore satisfies the hypotheses of
Theorem~\ref{thm:scar-blackbox}, and
\[
    \log_4\frac{a_x}{\mu}
    \le
    \log_4\kappa.
\]
Substitution into \eqref{eq:scar-original-complexity} proves
\eqref{eq:scar-inner-complexity}.

When $\kappa=1$, strong monotonicity and Lipschitz continuity with the
same modulus imply the corresponding endpoint estimate directly, so
the cost of a residual-reduction call is bounded by a numerical
constant.
\end{proof}


\begin{proof}[Proof of Lemma~\ref{lem:pf-asymptotic-oracle}]
Let
\[
    s:=\sqrt{M_P/6}.
\]
Choose
\begin{equation}
    \epsilon_{\rm or}
    :=
    \min\left\{
        1,\,
        \frac{\mu}{\sqrt\rho},\,
        \left(
            \frac{\mu^2}{
                128(1+\kappa)\rho
            }
        \right)^{1/3},\,
        \left(
            \frac{
                s\mu
            }{
                64(1+\kappa)^2\rho
            }
        \right)^{2/3},\,
        \left(
            \frac{
                3\mu^3
            }{
                7\rho
            }
        \right)^{1/3}
    \right\}.
    \label{eq:pf-uniform-oracle-onset}
\end{equation}
By Lemma~\ref{lem:inner-residual-certificate},
\[
    \|y-y^*(x)\|
    \le
    \frac{\epsilon^2}{\mu}.
\]
The choice \eqref{eq:pf-uniform-oracle-onset} ensures
\[
    \frac{\epsilon^2}{\mu}
    \le
    \frac{\mu}{\rho}.
\]
Lemma~\ref{lem:corrected-errors} therefore gives
\begin{align*}
    \|\widehat g-\nabla P\|
    &\le
    \frac{(1+\kappa)\rho}{2\mu^2}\epsilon^4,
    \\
    \|\widehat H-\nabla^2P\|
    &\le
    \frac{(1+\kappa)^2\rho}{\mu}\epsilon^2,
    \\
    |\widehat P-P|
    &\le
    \frac{7\rho}{24\mu^3}\epsilon^6.
\end{align*}
The remaining conditions in
\eqref{eq:pf-uniform-oracle-onset} make these quantities at most
\[
    \frac{\epsilon}{256},
    \qquad
    \frac{s\sqrt\epsilon}{64},
    \qquad
    \frac{\epsilon^3}{8},
\]
respectively.
All coefficients depend only on $(\ell,\mu,\rho)$, so the same
threshold applies uniformly over the entire parameter class.
\end{proof}


\begin{lemma}[Safe-scale robust trial]
\label{lem:pf-safe-trial}
Consider a nonzero trial of
Algorithm~\ref{alg:pf-residual-utr}, and write
\[
    h:=\max\{\|g\|,\epsilon\},
    \qquad
    r:=\|d\|.
\]
Suppose both current and candidate inner points satisfy
$\|\nabla_yf\|\le\epsilon^2$, and suppose
\eqref{eq:pf-oracle-bounds} holds.
If
\[
    \sigma^2\ge M_P/6,
\]
then
\begin{subequations}
\begin{align}
    P(x)-P(x+d)
    &\ge
    \frac14\sigma\sqrt h\,r^2
    +
    \frac12\lambda r^2
    -
    \frac{(1+\kappa)\rho}{2\mu^2}
    \epsilon^4r,
    \label{eq:pf-safe-decrease}
    \\
    \|g^+\|
    &\le
    \frac{117}{256}h+\lambda r.
    \label{eq:pf-safe-gradient}
\end{align}
\end{subequations}
If in addition
\[
    r=\frac{\sqrt h}{4\sigma},
\]
then
\begin{equation}
    P(x)-P(x+d)
    \ge
    \frac{h^{3/2}}{128\sigma}
    +
    \frac14\lambda r^2.
    \label{eq:pf-safe-boundary-decrease}
\end{equation}
\end{lemma}

\begin{proof}
Apply Lemma~\ref{lem:common-inexact-tr-step} with
\[
    F=P,
    \qquad
    M=M_P,
\]
and
\[
    a
    =
    \sigma\sqrt h
    +
    \frac{\sigma\sqrt\epsilon}{64},
    \qquad
    R
    =
    \frac{\sqrt h}{4\sigma},
\]
\[
    \xi
    =
    \frac1{64}
    \sqrt{\frac{M_P\epsilon}{6}},
    \qquad
    \delta_0=\delta_+
    =
    \frac{(1+\kappa)\rho}{2\mu^2}\epsilon^4.
\]
At a safe scale,
\[
    \xi
    \le
    \frac{\sigma\sqrt\epsilon}{64},
\]
so
\[
    a-\xi\ge\sigma\sqrt h,
    \qquad
    a+\xi
    \le
    \sigma\sqrt h
    +
    \frac{2\sigma\sqrt\epsilon}{64}.
\]
The trust-region radius and
$M_P\le6\sigma^2$ imply
\[
    \frac{M_P}{6}r^3
    \le
    \frac14\sigma\sqrt h\,r^2.
\]
Equation~\eqref{eq:common-function-bound} therefore proves
\eqref{eq:pf-safe-decrease}.

If the radius is active,
\[
    \sigma\sqrt h\,r^2
    =
    \frac{h^{3/2}}{16\sigma}.
\]
By \eqref{eq:pf-gradient-error}, the last term in
\eqref{eq:pf-safe-decrease} is at most
$h^{3/2}/(1024\sigma)$, which yields
\eqref{eq:pf-safe-boundary-decrease}.

Finally, \eqref{eq:common-gradient-bound} gives
\[
    \|g^+\|
    \le
    \sigma\sqrt h\,r
    +
    \lambda r
    +
    2\frac{\sigma\sqrt\epsilon}{64}r
    +
    \frac{(1+\kappa)\rho}{\mu^2}\epsilon^4
    +
    \frac{M_P}{2}r^2.
\]
The four terms other than $\lambda r$ are bounded by
\[
    \frac h4,
    \qquad
    \frac h{128},
    \qquad
    \frac h{128},
    \qquad
    \frac{3h}{16},
\]
respectively. Their sum is smaller than $117h/256$, which proves
\eqref{eq:pf-safe-gradient}.
\end{proof}


\begin{lemma}[Accepted-step accounting]
\label{lem:pf-accepted-accounting}
Consider any finite prefix of $N$ accepted nonterminal steps of
Algorithm~\ref{alg:pf-residual-utr}.
Suppose \eqref{eq:pf-oracle-bounds} holds at every stored tuple and
all trial scales are bounded by $\bar\sigma$.
If
\begin{equation}
    \frac{
        5\bar\sigma\epsilon^{3/2}
    }{
        4\eta
    }
    \left(
        1+\frac4{\log2}
    \right)
    \le
    \frac12,
    \label{eq:pf-absorption-condition}
\end{equation}
then
\begin{equation}
    N
    =
    O\!\left(
        \bar\sigma\Delta_P\epsilon^{-3/2}
        +
        \log_+\frac{\|g_0\|}{\epsilon}
        +1
    \right).
    \label{eq:pf-accepted-bound}
\end{equation}
\end{lemma}

\begin{proof}
For an accepted step $k$, write
\[
    h_k:=\max\{\|g_k\|,\epsilon\},
    \qquad
    r_k:=\|d_k\|,
\]
and define
\[
    \mathcal B_N
    :=
    \left\{
        0\le k<N:
        r_k=\frac{\sqrt{h_k}}{4\sigma_k}
    \right\},
\]
\[
    \mathcal I_N
    :=
    \left\{
        0\le k<N:
        r_k<\frac{\sqrt{h_k}}{4\sigma_k}
    \right\}.
\]
Every $k\in\mathcal B_N$ gives a true decrease of at least
\[
    \eta
    \left(
        \frac{h_k^{3/2}}{\sigma_k}
        +
        \lambda_kr_k^2
    \right),
\]
because the two corrected-value errors sum to at most
$\epsilon^3/4$.
Every $k\in\mathcal I_N$ can increase the true primal value by at
most $5\epsilon^3/4$.
Telescoping and using
$h_k\ge\epsilon$ and $\sigma_k\le\bar\sigma$ give
\begin{subequations}
\label{eq:pf-boundary-budgets}
\begin{align}
    |\mathcal B_N|
    &\le
    \frac{
        \bar\sigma\Delta_P
    }{
        \eta\epsilon^{3/2}
    }
    +
    \frac{
        5\bar\sigma\epsilon^{3/2}
    }{
        4\eta
    }
    |\mathcal I_N|,
    \label{eq:pf-boundary-count}
    \\
    \sum_{k\in\mathcal B_N}
    \frac{\lambda_kr_k}{h_k}
    &\le
    \frac{
        4\bar\sigma\Delta_P
    }{
        \eta\epsilon^{3/2}
    }
    +
    \frac{
        5\bar\sigma\epsilon^{3/2}
    }{
        \eta
    }
    |\mathcal I_N|.
    \label{eq:pf-multiplier-budget}
\end{align}
\end{subequations}
For the second inequality we used
\[
    \lambda_kr_k^2
    =
    \frac{\lambda_kr_k}{h_k}
    \frac{h_k^{3/2}}{4\sigma_k}.
\]

The accepted gradient tests imply
\[
    h_{k+1}
    \le
    \begin{cases}
        \max\{h_k/2,\epsilon\},
        &k\in\mathcal I_N,
        \\
        \max\{h_k/2+\lambda_kr_k,\epsilon\},
        &k\in\mathcal B_N.
    \end{cases}
\]
Thus every interior step that does not hit the floor decreases
$\log(h_k/\epsilon)$ by at least $\log2$, while a boundary step can
increase it by at most $\lambda_kr_k/h_k$.

Suppose an interior step hits the floor and $k+1<N$.
At the next accepted outer iteration, $h=\epsilon$ and the stopping
test fails through its Hessian branch.
With
\[
    A
    :=
    H
    +
    \frac{\sigma\sqrt\epsilon}{64}I
    +
    \sigma\sqrt\epsilon\,I,
\]
failure of the stopping test gives
\[
    -\lambda_{\min}(A)
    >
    \frac{7}{64}\sigma\sqrt\epsilon.
\]
Lemma~\ref{lem:tr-global-optimality} therefore forces
$\lambda>0$, and hence the next accepted step lies on the boundary.
Thus every nonterminal floor event is injectively followed by an
index in $\mathcal B_N$, apart from a possible final prefix index.
Therefore
\begin{equation}
    |\mathcal I_N|
    \le
    |\mathcal B_N|+1
    +
    \frac{
        \log_+(\|g_0\|/\epsilon)
        +
        \sum_{k\in\mathcal B_N}
        \lambda_kr_k/h_k
    }{
        \log2
    }.
    \label{eq:pf-interior-potential}
\end{equation}
Substituting \eqref{eq:pf-boundary-budgets} and using
\eqref{eq:pf-absorption-condition} proves
\eqref{eq:pf-accepted-bound}.
\end{proof}


\begin{proof}[Proof of Theorem~\ref{thm:pf-scar-utr}]
Set
\[
    s:=\sqrt{M_P/6},
    \qquad
    \bar\sigma:=2s.
\]
Choose the class-uniform threshold no larger than
$\epsilon_{\rm or}$ and so that
\begin{equation}
\begin{gathered}
    \epsilon\le s,
    \\
    \frac{
        5\bar\sigma\epsilon^3
    }{
        4\eta\epsilon^{3/2}
    }
    \left(
        1+\frac4{\log2}
    \right)
    \le
    \frac12,
    \\
    \frac1{s\sqrt\epsilon}
    \left(
        \frac{(1+\kappa)\rho}{2\mu^2}
        \epsilon^4
    \right)^2
    \le
    \frac34\epsilon^3,
    \\
    \left(
        \frac1{128}-\eta
    \right)
    \frac{\epsilon^{3/2}}{\bar\sigma}
    \ge
    \frac54\epsilon^3.
\end{gathered}
\label{eq:pf-small-epsilon-conditions}
\end{equation}
For example, the choice
\begin{equation}
    \epsilon_0
    :=
    \frac12
    \min\left\{
        \epsilon_{\rm or},
        s,
        \left(
            \frac{
                \eta
            }{
                5s(1+4/\log2)
            }
        \right)^{2/3},
        \left(
            \frac{
                3s\mu^4
            }{
                (1+\kappa)^2\rho^2
            }
        \right)^{2/9}
    \right\}
    \label{eq:pf-class-uniform-onset}
\end{equation}
is sufficient.
This threshold depends only on $(\ell,\mu,\rho)$.

Starting from $\sigma=\epsilon$, a rejected trial doubles $\sigma$.
Unless the method has already terminated, the first nonzero trial
with $\sigma\ge s$ satisfies
\[
    \sigma\le2s=\bar\sigma.
\]
Lemma~\ref{lem:pf-safe-trial} implies that its gradient acceptance
test holds.

For an interior trial, minimizing the right-hand side of
\eqref{eq:pf-safe-decrease} over $r\ge0$ shows that the possible true
increase is at most
\[
    \frac1{\sigma\sqrt h}
    \left(
        \frac{(1+\kappa)\rho}{2\mu^2}
        \epsilon^4
    \right)^2.
\]
By \eqref{eq:pf-small-epsilon-conditions}, this is at most
$3\epsilon^3/4$. After adding the two corrected-value errors, the
interior value test holds.

On a boundary trial,
\eqref{eq:pf-safe-boundary-decrease} and the two value errors give
\[
    \widehat P-\widehat P^+
    \ge
    \frac{h^{3/2}}{128\sigma}
    +
    \frac14\lambda r^2
    -
    \frac{\epsilon^3}{4}.
\]
The last condition in
\eqref{eq:pf-small-epsilon-conditions} implies the boundary
acceptance test.
Hence a safe-scale trial either terminates or is accepted, and no
generated scale exceeds $\bar\sigma$.

Consider a finite trial prefix with $K_N$ accepted and $U_N$
rejected nonterminal trials.
Lemma~\ref{lem:pf-accepted-accounting} bounds $K_N$ independently of
the prefix length.
Because each rejection doubles $\sigma$ and each acceptance decreases
it by at most a factor two,
\[
    U_N
    \le
    K_N+
    \log_2(\bar\sigma/\epsilon)+1.
\]
Hence every finite prefix has a common finite length bound, so an
infinite run is impossible.

For the complete run,
\[
    Q
    \le
    K+U+1.
\]
Moreover,
\[
    \|g_0\|
    \le
    \|\nabla P(x_0)\|
    +
    \frac{\epsilon}{256},
\]
so the initial logarithm can be replaced by
$\log_+(\|\nabla P(x_0)\|/\epsilon)$ up to a universal additive
constant. This proves \eqref{eq:pf-outer-complexity}.

If the explicit stopping test returns, then
\eqref{eq:pf-gradient-error},
\eqref{eq:pf-hessian-error},
and $\sigma\le2s$ give
\[
    \|\nabla P(x)\|
    \le
    \frac{257}{256}\epsilon,
\]
and
\[
    \lambda_{\min}(\nabla^2P(x))
    \ge
    -\frac{145}{64}s\sqrt\epsilon
    \ge
    -\sqrt{M_P\epsilon}.
\]
If the robust subproblem returns $d=0$, global optimality gives
$g=0$, $h=\epsilon$, and
\[
    H+
    \frac{\sigma\sqrt\epsilon}{64}I
    +
    \sigma\sqrt\epsilon I
    \succeq0.
\]
The same oracle bounds then imply
\[
    \|\nabla P(x)\|
    \le
    \frac{\epsilon}{256},
    \qquad
    \lambda_{\min}(\nabla^2P(x))
    \ge
    -\frac{131}{64}s\sqrt\epsilon
    \ge
    -\sqrt{M_P\epsilon}.
\]
This proves \eqref{eq:pf-return}.
\end{proof}


\begin{proof}[Proof of Corollary~\ref{cor:pf-scar-inner-complexity}]
Every accepted boundary step decreases the true primal value, while
every accepted interior step can increase it by at most
$5\epsilon^3/4$.
Since there are at most $Q$ accepted steps,
\[
    P(x)-P^*
    \le
    \Delta_P+\frac54Q\epsilon^3
\]
at every accepted outer point.

Because $\nabla P$ is $(1+\kappa)\ell$-Lipschitz, the standard
descent inequality together with
\eqref{eq:pf-gradient-error} gives
\[
    h\le G_Q.
\]

At a candidate point, SCAR is warm-started from a point whose previous
inner residual is at most $\epsilon^2$.
Joint $\ell$-smoothness and the trust-region radius give
\[
    \|\nabla_yf(x+d,y)\|
    \le
    \epsilon^2+\ell\|d\|
    \le
    \epsilon^2
    +
    \frac{\ell\sqrt{G_Q}}{4\epsilon}.
\]
Apply Corollary~\ref{cor:scar-inner} with target $\epsilon^2$ to
every candidate call and separately to the initial call.
There is at most one candidate call per solved subproblem, and
summing the resulting bounds gives
\eqref{eq:pf-inner-complexity}.
\end{proof}


\subsection{Persistent Inner Tracking}
\label{app:persistent-scar-proofs}


\begin{proof}[Proof of Lemma~\ref{lem:scar-persistent-half}]
The AR estimate used in the proof of
\citet[Theorem~5.1]{lan2026optimal}, obtained from their
Theorem~4.1, states that a call with current guess $\nu$ and an
admissible line-search estimate uses at most
\[
    4+C_1\sqrt{10L/\nu}
\]
gradient evaluations and returns $v$ satisfying
\[
    \|\nabla\phi_j(v)\|
    \le
    \frac{\nu}{2}
    \|u-u_j^*\|,
    \qquad
    u_j^*
    :=
    \arg\min_u\phi_j(u).
\]
If $\nu\le\mu$, strong convexity gives
\[
    \|u-u_j^*\|
    \le
    \frac{\|\nabla\phi_j(u)\|}{\mu},
\]
so the returned residual is at most $R/2$.
Therefore a failed halving implies $\nu>\mu$.

Every failure divides $\nu$ by four.
Once $\nu\le\mu$, no later objective in the sequence can cause
another curvature-guess failure.
Thus the total number of failures is at most
\[
    \left\lceil
        \log_4\frac{\nu_0}{\mu}
    \right\rceil.
\]
Throughout the process,
\[
    \nu\ge\mu/4,
\]
so every attempted stage costs at most
\[
    4+C_1\sqrt{40\kappa}
    \le
    (4+2\sqrt{10}\,C_1)\sqrt\kappa.
\]
There are $N$ successful attempts plus at most the displayed number
of failed attempts, which proves
\eqref{eq:persistent-scar-count}.
The line-search estimate follows from the corresponding SCAR bound.
\end{proof}


\begin{proof}[Proof of Lemma~\ref{lem:gs-tracking}]
Joint Lipschitz continuity of the gradient gives
\[
    \|\nabla_yf(x+d,y)\|
    \le
    \|\nabla_yf(x,y)\|
    +
    \ell r
    \le
    \ell
    \bigl(
        c_0r_\epsilon(\sigma)+r
    \bigr).
\]
After
$N_{\rm tr}(r,\sigma)$ successful residual halvings, the right-hand
side is at most
\[
    c_0\ell r_\epsilon(\sigma),
\]
which proves the tracking claim.
At an unchanged $x$,
\[
    r_\epsilon(2\sigma)
    =
    \frac12r_\epsilon(\sigma),
\]
so one successful halving restores the invariant after scale
doubling.
\end{proof}


\begin{proof}[Proof of Lemma~\ref{lem:gs-working-errors}]
From strong concavity and
\eqref{eq:gs-working-invariant},
\[
    e
    :=
    \|y-y^*(x)\|
    \le
    \frac{
        c_0\kappa\sqrt\epsilon
    }{
        4\sigma
    }
    \le
    \frac{c_0\kappa}{4}
    \sqrt\epsilon\,
    \mathcal A_\epsilon.
\]
For all sufficiently small $\epsilon$, uniformly over the fixed
parameter class,
\[
    e\le\frac{\mu}{\rho}.
\]
Lemma~\ref{lem:corrected-errors} therefore gives
\[
    \|\widehat g-\nabla P(x)\|
    \le
    \frac{(1+\kappa)\rho}{2}e^2,
\]
\[
    |\widehat P-P(x)|
    \le
    \frac{7\rho}{24}e^3,
\]
and
\[
    \|\widehat H-\nabla^2P(x)\|
    \le
    (1+\kappa)^2\rho e.
\]
Substitution gives
\eqref{eq:gs-working-errors}.

At a safe scale,
\[
    M_P\le6\sigma^2.
\]
Using
$M_P=4\sqrt2\,\kappa^3\rho$ and $\kappa\ge1$ gives
\begin{align*}
    \|\widehat g-\nabla P(x)\|
    &\le
    \frac{3c_0^2}{32\sqrt2}\epsilon,
    \\
    \|\widehat H-\nabla^2P(x)\|
    &\le
    \frac{3\sqrt2c_0}{4}
    \sigma\sqrt\epsilon,
    \\
    |\widehat P-P(x)|
    &\le
    \frac{7c_0^3}{1024\sqrt2}
    \frac{\epsilon^{3/2}}{\sigma}.
\end{align*}
Since $c_0=2^{-12}$, these imply
\eqref{eq:working-safe-errors}.

For a validation point,
$\mu\le a_x\le\ell$ and
Lemma~\ref{lem:inner-residual-certificate} imply
\[
    \|y_{\rm v}-y^*(x)\|
    \le
    c_0\kappa\epsilon^{3/2}.
\]
The corrected-gradient and Hessian errors are therefore at most
\[
    \frac{
        (1+\kappa)\rho
        c_0^2\kappa^2
    }2
    \epsilon^3
\]
and
\[
    (1+\kappa)^2\rho
    c_0\kappa
    \epsilon^{3/2},
\]
respectively.
The first is at most $\epsilon/1024$ for sufficiently small
$\epsilon$.
Since
\[
    \epsilon\mathcal A_\epsilon\to0
\]
and
$\sigma\ge\mathcal A_\epsilon^{-1}$,
the second is at most
$\sigma\sqrt\epsilon/32$ for sufficiently small $\epsilon$.
This proves \eqref{eq:validator-errors}.
\end{proof}


\begin{proof}[Proof of Lemma~\ref{lem:gs-safe}]
Apply Lemma~\ref{lem:common-inexact-tr-step} with
\[
    F=P,
    \qquad
    M=M_P,
\]
\[
    a
    =
    \sigma\sqrt h
    +
    \frac{\sigma\sqrt\epsilon}{64},
    \qquad
    \xi
    =
    \frac{\sigma\sqrt\epsilon}{64},
\]
and
\[
    \delta_0=\delta_+
    =
    \frac{\epsilon}{1024}.
\]
The endpoint errors follow from
\eqref{eq:working-safe-errors}.
Because
\[
    r\le\frac{\sqrt h}{4\sigma},
    \qquad
    M_P\le6\sigma^2,
\]
we have
\[
    \frac{M_P}{6}r^3
    \le
    \frac W4.
\]
Thus
\[
    P(x)-P(x+d)
    \ge
    \frac W4
    +
    \frac U2
    -
    \frac{\epsilon r}{1024}.
\]
Young's inequality and $h\ge\epsilon$ give
\[
    \frac{\epsilon r}{1024}
    \le
    \frac W8
    +
    \frac{2T}{1024^2}.
\]
After paying the two corrected-value errors,
\[
    \widehat P-\widehat P^+
    \ge
    \frac W8
    +
    \frac U2
    -
    \left(
        \frac2{1024^2}
        +
        \frac2{4096}
    \right)T
    \ge
    \frac W8
    +
    \frac U4
    -
    \frac T{1024}.
\]
The candidate-gradient estimate similarly gives
\[
    \|\widehat g^+\|
    \le
    \lambda r
    +
    \frac h4
    +
    \frac h{128}
    +
    \frac{3h}{16}
    +
    \frac h{512}
    =
    \lambda r+\frac{229}{512}h
    \le
    \lambda r+\frac h2.
\]
Hence every ordinary trial is accepted at a safe scale.

If $h=\epsilon$ and $r<R(h,\sigma)$, then complementarity gives
$\lambda=0$.
The candidate-gradient estimate yields
\[
    \|\nabla P(x+d)\|
    \le
    \frac{457}{1024}\epsilon.
\]
Positive semidefiniteness of the trust-region model, the working
Hessian error, and Hessian Lipschitz continuity give
\[
    \lambda_{\min}
    (\nabla^2P(x+d))
    \ge
    -\frac{81}{32}
    \sigma\sqrt\epsilon.
\]
Using the validation errors,
\[
    \|\widehat g_{\rm v}\|
    \le
    \frac{458}{1024}\epsilon,
\]
and
\[
    \lambda_{\min}(\widehat H_{\rm v})
    \ge
    -\frac{82}{32}
    \sigma\sqrt\epsilon,
\]
so validation succeeds.

Finally, successful validation at any generated scale gives
\[
    \|\nabla P(x^+)\|
    <
    \epsilon
\]
and
\[
    \lambda_{\min}(\nabla^2P(x^+))
    \ge
    -
    \left(
        \frac{21}{8}
        +
        \frac1{32}
    \right)
    \sigma\sqrt\epsilon
    =
    -\frac{85}{32}
    \sigma\sqrt\epsilon.
\]
\end{proof}


\begin{lemma}[Accepted-step potential decrease]
\label{lem:gs-descent}
Every ordinary accepted step satisfies
\begin{equation}
    \Phi-\Phi^+
    \ge
    \frac W{32}
    +
    \frac T{1024}
    +
    \frac U{16}.
    \label{eq:gs-potential-decrease-app}
\end{equation}
\end{lemma}

\begin{proof}
If $h>\epsilon$ and
\[
    h^+\le\max\{h/2,\epsilon\},
\]
then
\begin{equation}
    \mathcal L_\epsilon(h^+)
    \le
    \mathcal L_\epsilon(h)-1.
    \label{eq:gs-layer-halving}
\end{equation}
More generally, if $c\ge1$ and $h^+\le ch$, then
\begin{equation}
    \mathcal L_\epsilon(h^+)
    -
    \mathcal L_\epsilon(h)
    \le
    \lceil\log_2c\rceil
    \le
    1+\log_2c.
    \label{eq:gs-layer-growth}
\end{equation}

For an interior accepted step,
$\lambda=0$ and $U=0$.
The gradient acceptance test and
\eqref{eq:gs-layer-halving} give
\[
    \Phi-\Phi^+
    \ge
    \frac W8
    -
    \frac T{1024}
    +
    \frac T{512}
    =
    \frac W8+\frac T{1024}.
\]

For a boundary step, let
\[
    z:=\frac{\lambda r}{h}.
\]
The gradient acceptance test gives
\[
    \frac{h^+}{h}
    \le
    \max\{1,1/2+z\}.
\]
Hence
\[
    \mathcal L_\epsilon(h^+)
    -
    \mathcal L_\epsilon(h)
    \le
    1+\frac z{\ln2}.
\]
Since
\[
    r=\frac{\sqrt h}{4\sigma},
\]
we have
\[
    W
    =
    \frac{h^{3/2}}{16\sigma}
    \ge
    \frac T{16},
\]
\[
    U
    =
    \frac{zh^{3/2}}{4\sigma},
    \qquad
    Tz\le4U.
\]
Therefore
\begin{align*}
    \Phi-\Phi^+
    &\ge
    \frac W8
    +
    \frac U4
    -
    \frac T{1024}
    -
    \frac T{512}
    \left(
        1+\frac z{\ln2}
    \right)
    \\
    &\ge
    \frac{5W}{64}
    +
    \left(
        \frac14-\frac1{128\ln2}
    \right)U.
\end{align*}
Finally,
\[
    \frac{5W}{64}
    \ge
    \frac W{32}
    +
    \frac T{1024},
\]
and
\[
    \frac14-\frac1{128\ln2}
    \ge
    \frac1{16},
\]
which proves the claim.
\end{proof}


Define
\begin{equation}
\begin{aligned}
    E(\sigma)
    &:=
    C_v\frac{\epsilon^{3/2}}{\sigma^3},
    &
    E_0
    &:=
    E(\sigma_0),
    \\
    T_0
    &:=
    \frac{\epsilon^{3/2}}{\sigma_0},
    &
    h_0
    &:=
    \max\{\|\widehat g_0\|,\epsilon\},
    \\
    \mathcal L_0
    &:=
    \mathcal L_\epsilon(h_0),
    &
    D_0
    &:=
    1+
    \log_2
    \left(
        1+\frac{5C_g}{4\sigma_0^2}
    \right).
\end{aligned}
\label{eq:gs-budget-parameters}
\end{equation}
Set
\begin{equation}
    B_\epsilon
    :=
    \Delta_P
    +
    \frac{23}{7}E_0
    +
    \frac{T_0}{512}
    (\mathcal L_0+D_0).
    \label{eq:gs-budget}
\end{equation}

For sufficiently small $\epsilon$, let
\[
    \sigma_*:=\sqrt{M_P/6}.
\]
Lemma~\ref{lem:gs-safe} implies that every failed trial occurs below
$\sigma_*$.
Hence, for every finite prefix,
\begin{equation}
    J
    \le
    \left\lceil
        \log_2(
            \sigma_*\mathcal A_\epsilon
        )
    \right\rceil,
    \qquad
    \sigma
    \le
    \overline\sigma
    :=
    2\sigma_*.
    \label{eq:gs-scale-bound}
\end{equation}

\begin{lemma}[Cumulative accepted-step budget]
\label{lem:gs-budget}
Every finite prefix satisfies
\begin{equation}
    \sum_{\rm accepted}
    \left(
        \frac{W_k}{32}
        +
        \frac{T_k}{1024}
        +
        \frac{U_k}{16}
    \right)
    \le
    B_\epsilon.
    \label{eq:gs-cumulative-budget}
\end{equation}
If $K$ denotes its number of accepted steps, then
\begin{equation}
    K
    \le
    1024\overline\sigma
    B_\epsilon
    \epsilon^{-3/2},
    \qquad
    \sum_{\rm accepted}
    \sigma_k^2r_k^3
    \le
    8B_\epsilon.
    \label{eq:gs-path-budget}
\end{equation}
Every stored iterate satisfies
\[
    P(x)-P^*
    \le
    B_\epsilon,
\]
and
\begin{equation}
    B_\epsilon
    =
    \Delta_P
    +
    O_{\rm fixed}
    \bigl(
        \epsilon^{3/2}\mathcal A_\epsilon^3
    \bigr).
    \label{eq:gs-budget-asymptotic-app}
\end{equation}
\end{lemma}

\begin{proof}
When the scale is refreshed from $\sigma$ to $2\sigma$ at an
unchanged outer point,
\eqref{eq:gs-working-errors} bounds the corrected-value jump by
\[
    E(\sigma)+E(2\sigma)
    =
    \frac98E(\sigma).
\]
The failure scales form a doubling sequence, so the sum of all
positive corrected-value jumps is at most
\[
    \frac98E_0
    \sum_{j\ge0}2^{-3j}
    =
    \frac97E_0.
\]

For the gradient-level contribution,
the two corrected-gradient errors imply
\[
    \|\widehat g_{\rm new}
    -
    \widehat g_{\rm old}\|
    \le
    \frac{5C_g\epsilon}{4\sigma^2}.
\]
Since
$g\mapsto\max\{\|g\|,\epsilon\}$ is $1$-Lipschitz,
\[
    \frac{h_{\rm new}}{h_{\rm old}}
    \le
    1+
    \frac{5C_g}{4\sigma^2},
\]
and therefore
\[
    \mathcal L_\epsilon(h_{\rm new})
    -
    \mathcal L_\epsilon(h_{\rm old})
    \le
    D_0.
\]
Because the corresponding
$T=\epsilon^{3/2}/\sigma$ is halved after scale doubling, the positive
jump of the level term is at most
\[
    \frac{T_{\rm old}D_0}{1024}.
\]
Summing the doubling sequence gives a total at most
\[
    \frac{T_0D_0}{512}.
\]

The initial potential is bounded above by
\[
    P(x_0)
    +
    E_0
    +
    \frac{T_0}{512}\mathcal L_0,
\]
while every stored potential is bounded below by
$P^*-E_0$.
Telescoping
Lemma~\ref{lem:gs-descent}
and paying the two types of positive scale-refresh jumps gives
\[
    \sum_{\rm accepted}
    \left(
        \frac W{32}
        +
        \frac T{1024}
        +
        \frac U{16}
    \right)
    \le
    \Delta_P
    +
    2E_0
    +
    \frac97E_0
    +
    \frac{T_0}{512}
    (\mathcal L_0+D_0)
    =
    B_\epsilon.
\]
This proves
\eqref{eq:gs-cumulative-budget}.

Since
\[
    T_k
    \ge
    \frac{\epsilon^{3/2}}{\overline\sigma},
\]
the first estimate in
\eqref{eq:gs-path-budget} follows.
The radius constraint gives
\[
    4\sigma_k^2r_k^3
    \le
    \sigma_k\sqrt{h_k}\,r_k^2
    =
    W_k,
\]
which proves the second estimate.

Finally,
\[
    E_0
    =
    C_v\epsilon^{3/2}
    \mathcal A_\epsilon^3,
    \qquad
    T_0
    =
    \epsilon^{3/2}
    \mathcal A_\epsilon.
\]
At the initial stored point,
\[
    \|\widehat g_0\|
    \le
    \|\nabla P(x_0)\|
    +
    C_g\epsilon
    \mathcal A_\epsilon^2.
\]
For every fixed payoff and initialization,
\[
    \mathcal L_0
    =
    O_{\rm fixed}(\mathcal A_\epsilon),
\]
and
\[
    D_0
    =
    O_{\rm fixed}
    (1+\log\mathcal A_\epsilon).
\]
Substitution into \eqref{eq:gs-budget} gives
\eqref{eq:gs-budget-asymptotic-app}.
\end{proof}


\begin{proof}[Proof of Theorem~\ref{thm:gradient-scaled-scar}]
Choose $\epsilon_0$ so that
Lemmas~\ref{lem:gs-working-errors} and~\ref{lem:gs-safe} hold and
\[
    \mathcal A_\epsilon^{-1}
    <
    \sqrt{M_P/6}
\]
for every $0<\epsilon\le\epsilon_0$.
These restrictions depend only on the common constants
$(\ell,\mu,\rho)$.

Every inner sequence is finite by
Lemma~\ref{lem:scar-persistent-half}.
Every nonterminal outer trial is either accepted or increases
$\sigma$.
The finite-prefix bounds
\eqref{eq:gs-scale-bound} and
\eqref{eq:gs-path-budget}
therefore imply finite termination.

Let $K$ and $J$ denote the numbers of accepted steps and failed
trials. Then
\[
    Q=K+J+1,
\]
and
\[
    Q
    \le
    1024\overline\sigma
    B_\epsilon
    \epsilon^{-3/2}
    +
    \left\lceil
        \log_2(
            \sigma_*
            \mathcal A_\epsilon
        )
    \right\rceil
    +1.
\]
Using
\eqref{eq:gs-budget-asymptotic-app} and
\[
    \overline\sigma
    =
    2\sqrt{M_P/6}
\]
gives
\eqref{eq:gs-outer-complexity}.
The output guarantee follows from
\eqref{eq:gs-validator-return} and the scale cap.

It remains to count successful residual halvings.
For an accepted step, set
\[
    z_k
    :=
    \frac{
        \sigma_kr_k
    }{
        \sqrt\epsilon
    }.
\]
Since $c_0=2^{-12}$,
\[
    N_{{\rm tr},k}
    =
    \left\lceil
        \log_2(1+2^{14}z_k)
    \right\rceil
    \le
    16+z_k^3.
\]
Therefore
\begin{align}
    \sum_{\rm accepted}
    N_{{\rm tr},k}
    &\le
    16K
    +
    \epsilon^{-3/2}
    \sum_{\rm accepted}
    \sigma_k^3r_k^3
    \notag\\
    &\le
    16K
    +
    8\overline\sigma
    B_\epsilon
    \epsilon^{-3/2}.
    \label{eq:gs-accepted-tracking}
\end{align}

To control failed trials, let
\[
    L_P:=(1+\kappa)\ell.
\]
Smoothness of $P$ and its lower bound give
\[
    \|\nabla P(x)\|^2
    \le
    2L_P(P(x)-P^*).
\]
Hence every stored iterate satisfies
\[
    h
    \le
    \overline h_\epsilon
    :=
    \max\left\{
        \epsilon,\,
        \sqrt{2L_PB_\epsilon}
        +
        C_g\epsilon
        \mathcal A_\epsilon^2
    \right\}
    =
    O_{\rm fixed}(1).
\]
Each failed trial therefore requires at most
\[
    N_{\rm fail}
    :=
    \left\lceil
        \log_2\left(
            1+
            \frac1{c_0}
            \sqrt{
                \frac{\overline h_\epsilon}{\epsilon}
            }
        \right)
    \right\rceil
    =
    O_{\rm fixed}(\mathcal A_\epsilon)
\]
tracking halvings.

A terminal working update uses at most
\[
    \left\lceil
        \log_2(1+1/c_0)
    \right\rceil
    =
    13
\]
halvings, and each failure uses one additional halving for the scale
refresh.

Before validation, the working residual is at most
\[
    \frac{
        c_0\ell\sqrt\epsilon
    }{
        4\sigma
    },
\]
while the validation target is at least
\[
    c_0\mu\epsilon^{3/2}.
\]
Thus every validation requires at most
\[
    N_{\rm val}
    :=
    \left\lceil
        \log_2
        \max\left\{
            1,\,
            \frac{
                \kappa\mathcal A_\epsilon
            }{
                4\epsilon
            }
        \right\}
    \right\rceil
    =
    O_{\rm fixed}(\mathcal A_\epsilon)
\]
additional successful halvings.
Initialization similarly uses
\[
    N_{\rm init}
    :=
    \left\lceil
        \log_2
        \max\left\{
            1,\,
            \frac{
                \|\nabla_yf(x_0,y_{-1})\|
            }{
                c_0a_0r_\epsilon(\sigma_0)
            }
        \right\}
    \right\rceil
    =
    O_{\rm fixed}(\mathcal A_\epsilon)
\]
halvings.

Consequently,
\[
    N_{\rm half}
    \le
    \sum_{\rm accepted}N_{{\rm tr},k}
    +
    JN_{\rm fail}
    +
    J+13
    +
    (J+1)N_{\rm val}
    +
    N_{\rm init}.
\]
Since
\[
    J
    =
    O_{\rm fixed}
    (1+\log\mathcal A_\epsilon),
\]
we obtain
\[
    N_{\rm half}
    =
    O\!\left(
        \Delta_P\sqrt{M_P}\,
        \epsilon^{-3/2}
    \right)
    +
    O_{\rm fixed}
    (\mathcal A_\epsilon^3).
\]

All inner objectives $q_x=-f(x,\cdot)$ share the same
$(\mu,\ell)$.
Applying Lemma~\ref{lem:scar-persistent-half} once to the complete
sequence, including rejected candidates, scale refreshes, and
validators, gives
\begin{equation}
    N_y
    \le
    (4+2\sqrt{10}\,C_1)
    \sqrt\kappa
    \left(
        N_{\rm half}
        +
        \lceil\log_4\kappa\rceil
    \right)
    +
    2(J+2).
    \label{eq:gs-inner-explicit}
\end{equation}
This proves
\eqref{eq:gs-inner-complexity}.
Finally,
\[
    \sqrt{M_P}
    =
    \Theta(\kappa^{3/2}\sqrt\rho),
    \qquad
    \sqrt{\kappa M_P}
    =
    \Theta(\kappa^2\sqrt\rho),
\]
which gives \eqref{eq:gs-leading-complexities}.
\end{proof}